\section{HHL and Quantum Hoare Logic}\label{sec:hhl}

\subsection{The Algorithm}

Write $A=\sum_j\lambda_j\ketbra{u_j}{u_j}$ and $\ket b=\sum_j\beta_j\ket{u_j}$. Phase estimation with
the unitaries $e^{iAt_0}$ writes an approximation $\tilde\lambda_j$ of each eigenvalue into a clock
register, which gives $\sum_j\beta_j\ket{\tilde\lambda_j}\ket{u_j}$. A rotation of an ancilla qubit,
controlled by the clock, produces the amplitudes $C/\tilde\lambda_j$. Uncomputing the clock and
measuring the ancilla leaves, on outcome~$1$, the state $\sum_j(\beta_j/\lambda_j)\ket{u_j}$, which is
$A^{-1}\ket b$ up to normalisation. The success probability is $p=C^2\norm{A^{-1}b}^2$, at least
$1/\kappa^2$ for $C=\lambda_{\min}$ and unit~$b$, so the measurement is repeated or amplified.
Figure~\ref{fig:hhl} shows the circuit together with the state after each gate, which is itself a
theorem of the development.

\begin{figure}[t]
\centering
\scalebox{0.8}{%
\begin{quantikz}[row sep={0.8cm,between origins}, column sep=0.2cm]
\lstick{clock $\ket{0}^{\otimes k}$} & \qwbundle{k} &  & \gate{H^{\otimes k}}\slice{$\psi_2$}
  & \ctrl{1}\slice{$\psi_3$} & \gate{F^\dagger}\slice{$\psi_4$} & \ctrl{2}\slice{$\psi_5$}
  & \gate{F}\slice{$\psi_6$} & \ctrl{1}\slice{$\psi_7$} & \gate{H^{\otimes k}}\slice{$\psi_8$} & \\
\lstick{system $\ket{0}^{\otimes s}$} & \qwbundle{s} & \gate{U_b}\slice{$\psi_1$} & & \gate{U^{c}}
  & & & & \gate{U^{-c}} & & \rstick{$\propto A^{-1}b$ if $1$}\\
\lstick{ancilla $\ket0$} & & & & & & \gate{R(C/\lambda_c)} & & & & \meter{}
\end{quantikz}}
\\[0.6em]
\renewcommand{\arraystretch}{1.15}\small
\begin{tabular}{@{}l@{\enspace}l@{\qquad}l@{\enspace}l@{}}
$\psi_1$ & $\ket0\otimes b\otimes\ket0$ &
$\psi_5$ & $\sum_j\ket j\otimes\Pi_jb\otimes R(C/\lambda_j)\ket0$\\
$\psi_2$ & $u\otimes b\otimes\ket0$ &
$\psi_6$ & $\sum_j\ket j\otimes N^{-1/2}(U^{j}\otimes I)\,w$\\
$\psi_3$ & $\sum_j\ket j\otimes N^{-1/2}U^{j}b\otimes\ket0$ &
$\psi_7$ & $u\otimes w$\\
$\psi_4$ & $\sum_j\ket j\otimes\Pi_jb\otimes\ket0$ &
$\psi_8$ & $\ket0\otimes w$
\end{tabular}\\[0.2em]
{\small $w=\sqrt{I-C^2A^{-2}}\,b\otimes\ket0+CA^{-1}b\otimes\ket1$}
\caption{The HHL circuit and the state after each of its eight gates, from $\ket0\ket0\ket0$, for
$N=2^k$, $U=e^{iAt_0}$, $u=N^{-1/2}\sum_j\ket j$ and every eigenvalue of $A$ on the grid
$\lambda_j=2\pi j/(Nt_0)$ (read as $2\pi(j-N)/(Nt_0)$ for $j\ge N/2$). Gate~3 applies $U^c$ when the
clock holds $c$, gate~5 applies $R(C/\lambda_c)$ with $R(c)\ket0=\sqrt{1-c^2}\ket0+c\ket1$, and $\Pi_j$
is the spectral projector of $A$ at $\lambda_j$. The table is the theorem \isa{hhl\_steps}.}\label{fig:hhl}
\end{figure}

\subsection{Programs and Triples}

We use the concrete matrices of the Jordan-Normal-Form library~\cite{thiemann2015jnf}, shared by
\IMD~\cite{bordg2021imd,bordg2020afp}, the quantum Fourier transform entry~\cite{manrique2025qft} and
QHLProver~\cite{liu2019qhl,liu2019afp}. Dimensions are hypotheses (\isa{A $\in$ carrier\_mat n n})
rather than types, which is how the physics literature treats them and lets the libraries share
their definitions. QHLProver formalises Ying's quantum while-language~\cite{ying2011,ying2016book},
\[
S \;::=\; \mathbf{skip} \;\mid\; U \;\mid\; S_1;S_2 \;\mid\; \mathbf{measure}\ M\ [S_0,\dots,S_{m-1}]
\;\mid\; \mathbf{while}\ M\ \mathbf{do}\ S ,
\]
whose programs denote completely positive, trace-non-increasing maps on partial density operators;
the trace lost by a loop is the probability that it does not terminate. A predicate is an operator
$0\le P\le I$, and $\Tr(P\rho)$ is the probability that $P$ holds in the state~$\rho$. A triple
compares these probabilities before and after the program (Listing~\ref{lst:triples}).

\begin{lstlisting}[float=tp,caption={Triples in QHLProver (in the locale \isa{state\_sig} of the registers, of
total dimension \isa{d}; syntax annotations omitted)},label={lst:triples}]
definition hoare_total_correct :: "complex mat ~$\Rightarrow$~ com ~$\Rightarrow$~ complex mat
    ~$\Rightarrow$~ bool" where
  "~$\models_{t}$~ {P} S {Q} ~$\longleftrightarrow$~ (~$\forall\rho\in$~density_states. trace (P * ~$\rho$~) ~$\le$~ trace (Q *
      denote S ~$\rho$~))"
definition hoare_partial_correct :: "complex mat ~$\Rightarrow$~ com ~$\Rightarrow$~ complex
    mat ~$\Rightarrow$~ bool" where
  "~$\models_{p}$~ {P} S {Q} ~$\longleftrightarrow$~ (~$\forall\rho\in$~density_states.
      trace (P * ~$\rho$~) ~$\le$~ trace (Q * denote S ~$\rho$~) + (trace ~$\rho$~ - trace
          (denote S ~$\rho$~)))"
\end{lstlisting}

Every statement about HHL below is a triple of one of four shapes (Table~\ref{tab:readings}). Let
$P_0$, $P_1$ project onto ancilla value $0$, $1$, and let $\ket t$ be a unit vector with
$P_1\ket t=\ket t$. Since $\Tr\rho=1$ for a state, a precondition $q\cdot I$ asserts the number $q$ for
every input. The second reading needs a line of argument: if $\Tr((P_0+\ketbra tt)\sigma)\ge1$ for
$\sigma=\sem S\rho$, then $\Tr((P_1-\ketbra tt)\sigma)\le0$, and since $P_1-\ketbra tt\ge0$, the
ancilla can only read $1$ in the state $\ket t$.

\begin{table}[t]
\caption{The four shapes of triple used for HHL, for a state $\rho$ with $\Tr\rho=1$.}\label{tab:readings}
\centering\small
\begin{tabular}{@{}l>{\raggedright\arraybackslash}p{6.9cm}@{}}
\toprule
Triple & Reading\\
\midrule
$\hoaret{p\cdot I}{S}{P_1}$ & $S$ ends with the ancilla in $\ket1$ with probability at least $p$\\
$\hoaret{I}{S}{P_0+\ketbra tt}$ & whenever the ancilla reads $1$ at the end, the state is exactly $\ket t$\\
$\hoaret{I}{S}{\ketbra tt}$ & $S$ terminates with probability one, in the state $\ket t$\\
$\hoaret{q\cdot I}{S}{\ketbra tt}$ & the output has fidelity at least $q$ with $\ket t$\\
\bottomrule
\end{tabular}
\end{table}

Section~\ref{sec:programs} presents the proof system with which such triples are derived, the rules
we added to it for algorithms, and the triples of phase estimation and HHL.
